\begin{table}[htbp]
\centering
\caption{Mean-reversion diagnostic}
\label{tab:diagnostico_regressao_media}
\small
\begin{threeparttable}
\resizebox{\textwidth}{!}{%
\begin{tabular}{lcccc}
\toprule
Grade-subject & Main DR-DiD & Baseline support & $\Delta$ low controls & $\Delta$ high controls \\
\midrule
9th LP & 0.510 & 0.533 & 0.239 & -0.182 \\
9th Mat & 0.601 & 0.612 & -0.172 & -0.482 \\
HS3 LP & 0.571 & 0.581 & -0.094 & -0.578 \\
HS3 Mat & 0.627 & 0.607 & -0.295 & -0.542 \\
\bottomrule
\end{tabular}
}
\begin{tablenotes}
\footnotesize
\item \textit{Note:} The baseline-support column re-estimates DR-DiD restricting the sample to the p10--p90 interval of baseline proficiency among treated schools. The last two columns show the average change, in standard deviations, between 2011 and 2018 for never-treated controls below/above the treated-school baseline median.
\end{tablenotes}
\end{threeparttable}
\end{table}



